% !TeX program = pdflatex
\documentclass[aspectratio=169,11pt]{beamer}

\usetheme{Madrid}
\usecolortheme{default}
\usefonttheme{professionalfonts}
\setbeamercovered{transparent}

\setbeamertemplate{footline}{%
  \hfill\insertframenumber\,/\,\inserttotalframenumber\hspace{2mm}\vspace{2mm}%
}
\setbeamertemplate{navigation symbols}{}

\usepackage{amsmath,amssymb,mathtools,mathrsfs,bm,array}
\usepackage{bbm}
\usepackage{hyperref}

\definecolor{feiblue}{RGB}{22,78,128}
\definecolor{feigreen}{RGB}{45,110,62}
\definecolor{feired}{RGB}{145,48,48}
\definecolor{feilight}{RGB}{245,247,250}

\setbeamercolor{titlelike}{fg=white}
\setbeamercolor{frametitle}{fg=white,bg=feiblue}
\setbeamercolor{block title}{fg=white,bg=feiblue}
\setbeamercolor{block body}{fg=black,bg=feilight}
\setbeamercolor{alertblock title}{fg=white,bg=feired}
\setbeamercolor{alertblock body}{fg=black,bg=red!5}
\setbeamercolor{exampleblock title}{fg=white,bg=feigreen}
\setbeamercolor{exampleblock body}{fg=black,bg=green!5}

\AtBeginSection[]{
\begin{frame}[plain]
    \vfill
    \begin{center}
        {\usebeamerfont{title}\color{feiblue}\insertsectionhead}
    \end{center}
    \vfill
\end{frame}
}

\newcommand{\R}{\mathbb{R}}
\newcommand{\E}{\mathbb{E}}
\newcommand{\Pbb}{\mathbb{P}}
\newcommand{\cD}{\mathcal{D}}
\newcommand{\cP}{\mathcal{P}}
\newcommand{\cS}{\mathcal{S}}
\newcommand{\cX}{\mathcal{X}}
\newcommand{\cY}{\mathcal{Y}}
\newcommand{\dd}{\,\mathrm{d}}
\newcommand{\one}{\mathbbm{1}}
\newcommand{\norm}[1]{\left\lVert#1\right\rVert}
\newcommand{\ip}[2]{\left\langle#1,#2\right\rangle}
\newcommand{\onslidecite}[1]{\par\vspace{0.4ex}\hfill{\scriptsize\color{black!60}#1}}
\DeclareMathOperator{\KL}{KL}
\DeclareMathOperator{\softmax}{softmax}

\title{How Do Large Language Models Work?}
\subtitle{Mathematics of Scientific Foundation Models: motivation and overview}
\author{Fei Lu}
\date{Fall 2026 }

\begin{document}

%======================================================================
\begin{frame}
  \titlepage
\end{frame}

%======================================================================
\section{I. LLM: From Text to a Generated Sample}


%----------------------------------------------------------------------
\begin{frame}{One Short Prompt, Many Possible Continuations}
\begin{block}{Conditional law}
A short prompt $c$ rarely determines one continuation.  The mathematical
object should therefore be a conditional law
\[
  p_*(\,\cdot\mid c)\in\cP(\mathcal V).
\]
\end{block}

\begin{exampleblock}{Questions}
\begin{enumerate}
\item What conditional law is present in the data distribution?
\item How can one fixed Transformer evaluate a rich map $c\mapsto p_\theta(\cdot\mid c)$?
\item How does a decoding rule turn that law into one realized continuation?
\item How many conditional laws can a finite Transformer represent? ...
\end{enumerate}
\end{exampleblock}

\vspace{-0.5ex}
The empirical surprise of large-scale prompting is that changing $c$ can also
change the apparent task, while the model weights remain fixed.
\onslidecite{Brown et al. (2020)}
\end{frame}

\subsection{Token representation and next-token law}
%----------------------------------------------------------------------
\begin{frame}{}
\begin{center}
  {\Large  Text $\to$ Tokens $\to$ Causal Hidden States $\to$ Next-Token Law \\
  $\to$ Loss and Training $\to$ autoregressive generation}

  \vspace{6mm}
\begin{minipage}{0.92\textwidth}
\begin{itemize}
\item \textbf{Representation:} raw strings become finite token sequences.
\item \textbf{Computation:} causal attention maps a prefix to hidden states.
\item \textbf{Statistical target:} logits parameterize a next-token law.
\item \textbf{Learning:} cross-entropy fits that conditional law from data.
\item \textbf{Generation:} a decoding rule recursively samples or selects tokens.
\end{itemize}
\end{minipage}
\end{center}

\end{frame}

%----------------------------------------------------------------------
\begin{frame}{Tokenization Chooses the Finite State Space}
Let $\mathcal V$ be a finite vocabulary, $K=|\mathcal V|$.  A tokenizer is a
deterministic map
\[
 \operatorname{Tok}:\{\text{strings}\}
 \longrightarrow \bigcup_{T\ge 1}\mathcal V^T.
\]

\end{frame}

%----------------------------------------------------------------------
\begin{frame}{A Language Model Is a Product of Next-Token Conditionals}
Fix a horizon $T$ and let $U_{1:T} \in \mathcal V^T$ have probability distribution $p_*$. 
\begin{equation*}
 p_*(u_{1:T})
 =\prod_{t=1}^T p_*(u_t\mid u_{<t}),
 \qquad u_{<t}=(u_1,\ldots,u_{t-1}).
\end{equation*}

\begin{block}{Autoregressive model}
A parameter $\theta$ specifies
\[
 p_\theta(\,\cdot\mid u_{<t})\in\cP(\mathcal V),
 \qquad
 p_\theta(u_{1:T})
 =\prod_{t=1}^T p_\theta(u_t\mid u_{<t}).
\]
\end{block}
\begin{itemize}
  \item The learned object is a conditional probability distribution $p_\theta$.
  \item  Decoding, sampling, and search are subsequent decisions made from that vector.
\end{itemize}
\end{frame}



\subsection{Loss function, model architecture, and training}
%----------------------------------------------------------------------
\begin{frame}{The loss function}

To fit a model $p_\theta$ to data and estimate the population law $p_*$, define the population log-loss
\[
 \mathcal L(q)=\E_{p_*}\!\left[-\sum_{t=1}^T
     \log q(U_t\mid U_{<t})\right],
\]
\begin{itemize}
  \item The population law $p_*$ is unknown; use empirical approximation.
  \item It is the expected negative log-likelihood of the model $q$ under the true law $p_*$.
\end{itemize}

\end{frame}

%----------------------------------------------------------------------
\begin{frame}{Population Log-Loss and KL Divergence}
\begin{block}{Exact population identity}
For any admissible conditional law $q$,
\[
 \mathcal L(q)-\mathcal L(p_*)
 =\sum_{t=1}^T\E\!\left[
  \KL\!\left(p_*(\cdot\mid U_{<t})\,\middle\|\,
             q(\cdot\mid U_{<t})\right)
 \right]\ge 0.
\]
Equality holds iff the conditionals agree $p_*$-a.s. at every predicted position.
\end{block}
\begin{enumerate}
\item Condition on $U_{<t}$ and add and subtract
$-\log p_*(U_t\mid U_{<t})$.
\item Recall KL divergence: $\KL(p\|q)=\E_p[\log(p/q)]$.
\end{enumerate}


\begin{alertblock}{It does not prove}
that a finite Transformer contains $p_*$; that a finite data identifies it;
that a nonconvex optimizer finds the best parameter, or ....
\end{alertblock}
\end{frame}

%----------------------------------------------------------------------
\begin{frame}{Embeddings Carry Content, Position, Role}
For tokens $u_1,\ldots,u_N$, initialize
\[
 x_i^0=E_{\rm tok}e_{u_i}+r_i+s_i\in\R^d,
 \qquad i=1,\ldots,N.
\]
$N$ is typically less than or equal to $T$.  
\begin{center}
\begin{tabular}{c|c|c}
term & information & mathematical role \\
\hline
$E_{\rm tok}\in \R^{d\times K}$ & token embedding matrix &  \\
$e_{u_i}\in \R^K$ & token vector &  \\
$r_i$ & position & distinguishes order and distance \\
$s_i$ & role/modality & marks user, assistant, field, sensor, etc.
\end{tabular}
\end{center}

\begin{itemize}
  \item Masks, positions, and roles: break permutation symmetry in the attention map (later lectures). 

  
\end{itemize}

\end{frame}

%----------------------------------------------------------------------
\begin{frame}{One Causal Attention Head}
At layer $\ell$, let
\[
 q_i=Q\bar x_i^\ell,\qquad k_j=K\bar x_j^\ell,
 \qquad v_j=V\bar x_j^\ell.
\]
The masked weights and head output are
\begin{align*}
 a_{ij}
 &=\frac{\exp(\langle q_i,k_j\rangle/\sqrt{d_k})\,
            \textcolor{blue}{\one_{\{j\le i\}}}}
          {\sum_{m=1}^{\textcolor{blue}{i}}\exp(\langle q_i,k_m\rangle/\sqrt{d_k})},\\
 z_i&=\sum_{j=1}^i a_{ij}v_j.
\end{align*}

For each $i$, $a_{ij}\ge0$ and $\sum_{j\le i}a_{ij}=1$; hence $z_i$ lies in
the convex hull of the visible value vectors.

$$ \operatorname{Att}_\theta(X)_i=\sum_{j=1}^N a_{ij}Vx_j,\qquad a_{ij}\propto \exp(\langle Qx_i,Kx_j\rangle/\sqrt{d_k}).$$   
\onslidecite{Vaswani et al. (2017)}
\end{frame}


%----------------------------------------------------------------------
\begin{frame}{Residual Attention + Pointwise MLP = One Block}
For $H$ heads, a simplified pre-normalized residual block is
\begin{align*}
\text{Layer Normalization:} \quad  \bar x_i^\ell
 &=\operatorname{LN}_\ell(x_i^\ell), \quad z_i^{\ell h}= \mathrm{Attn}(X^{\ell}),  \, X = ( \bar x_i^\ell)_{i=1,\ldots,N} \\
  \text{Multi-Head Attention:}    \quad  \widetilde x_i^\ell
 &=x_i^\ell+W_O^\ell
   \operatorname{Concat}(z_i^{\ell1},\ldots,z_i^{\ell H}),\\
  x_i^{\ell+1}
 &=\widetilde x_i^\ell+W_2^\ell
   \sigma\!\left(W_1^\ell
     \operatorname{LN}'_\ell(\widetilde x_i^\ell)+b_1^\ell\right)+b_2^\ell.
\end{align*}
\vspace{-3mm}
After $L$ layers, $W_{\rm out}\in \R^{K\times d}, b_{\rm out}\in \R^K$ \vspace{-3mm} 
\[
 g_i=W_{\rm out}\operatorname{LN}_{\rm out}(x_i^L)+b_{\rm out},
 \qquad
 p_{\theta}(v\mid u_{\le i})=\frac{e^{(g_i)_v}}
 {\sum_{w\in\mathcal V}e^{(g_i)_w}}.
\]
\vspace{-3mm}
\begin{itemize}
  \item Attention (and multi-head) mixes positions; the MLP acts pointwise.
  \item Practical Transformers use different architectures, normalization, and activation functions. 
  \item $\theta = \{Q, K, V, W_O, W_1, W_2, b_1, b_2, W_{\rm out}, b_{\rm out}\}$ all weights in the model.
\end{itemize}

\end{frame}

%----------------------------------------------------------------------
\begin{frame}{ Training and Sequential Generation}

For a corpus $\cD_{\rm pre}=\{u_{1:T_r}^{(r)}\}_{r=1}^R$ with
$S=\sum_rT_r$, training minimizes
\[
 \widehat{\mathcal L}_S(\theta)
 =-\frac1S\sum_{r=1}^R\sum_{t=1}^{T_r}
 \log p_\theta\!\left(u_t^{(r)}\mid u_{<t}^{(r)}\right).
\]
\begin{itemize}
  \item All positions of a document are processed in parallel. 
  \item Parallel training and sequential generation are compatible. 
  \item Optimization is nonconvex and may not find the global minimum. ... 
\end{itemize}

\end{frame}



%----------------------------------------------------------------------
\begin{frame}{Softmax Gradient = Prediction $-$ Observation}
For logits $g\in\R^K$, target $y\in \{1,\ldots,K\}$, and $p=\softmax(g) = \frac{e^g}{\sum_j e^{g_j}}$,
\[
 \ell(g;y)= -\log p_y  
 =  -g_y+\log\sum_{j=1}^K e^{g_j}.
\]

\begin{block}{Exact local derivatives}
\[
 \nabla_g\ell(g;y)=p-e_y,
 \qquad
 \nabla_g^2\ell(g;y)=\operatorname{diag}(p)-pp^\top\succeq0.
\]
\[
 a^\top\nabla_g^2\ell\,a
 =\operatorname{Var}_{J\sim p}(a_J)\ge0, \quad \forall a\in\R^K.
\]
\end{block}
\begin{itemize}
  \item The gradient is the difference between the predicted probability and the observed class.
  \item The Hessian is a covariance matrix, hence positive semidefinite. EXE. 
  \item The token loss is convex in $g$, but not convex in the parameter $\theta$.
\end{itemize}

\end{frame}


%----------------------------------------------------------------------
\begin{frame}{Supervised Learning of Conditional Laws}
Each observed pair (prefix)
\[
  (U_{<t},U_t)
\]
is a conditional prediction problem, and one document supplies many such pairs.

\begin{block}{Empirical risk versus population risk}
\[
 \widehat{\mathcal L}_S(\theta)
 \quad\hbox{approximates}\quad
 \mathcal L(p_\theta)
 =\E_{p_*}\!\left[-\sum_{t=1}^T
     \log p_\theta(U_t\mid U_{<t})\right].
\]
\end{block}
\begin{itemize}
  \item  The prefixes within a document are dependent, not i.i.d. samples. 
  \item  $p_{\widehat \theta}$ aims to approximate $p_*$. 
  \item  The observed corpus need not be sample from one fixed law $p_*$. 
\end{itemize} 
Learning theory must specify the data process, model class, dependence, and
deployment distribution before assigning a generalization rate. 
\end{frame}

%----------------------------------------------------------------------
\begin{frame}{Autoregressive Sampling Closes the Loop}
Given a prompt $c=u_{1:N}$, repeat
\[
 U_{N+t}\sim p_\theta(\,\cdot\mid U_{1:N+t-1}).
\]
Then the generated continuation has exactly the product law
\[
 \Pbb_\theta(U_{N+1:N+T}=y_{1:T}\mid c)
 =\prod_{t=1}^T p_\theta(y_t\mid c,y_{<t}).
\]

\begin{itemize}
  \item The chain has closed: text $\to$ law $\to$ loss $\to$ a recursively generated sample. 
\item Post-processing, reranking, and feedback control, reinforcment learning are often used to improve quality. 
\end{itemize}

\end{frame}


%----------------------------------------------------------------------
\begin{frame}{Base Model, Instruction Tuning, Preference Tilt}
\begin{center}
pretraining on text
$\longrightarrow$
instruction tuning on prompt--response pairs
$\longrightarrow$
preference optimization
\end{center}

\begin{block}{Example: an idealized preference stage}
For base response law $p_\theta(\cdot\mid c)$ and reward $r(c,y)$,
\[
 \sup_{\pi(\cdot \mid c)}\left\{\E_\pi[r(c,Y)]
 -\beta\KL\bigl(\pi(\cdot\mid c)\|p_0(\cdot\mid c)\bigr)\right\}
\]
has optimizer
\[
 \pi^*(y\mid c)=
 \frac{p_\theta(y\mid c)e^{r(c,y)/\beta}}
      {\E_{p_\theta(\cdot\mid c)}[e^{r(c,Y)/\beta}]}.
\]
\end{block}

\onslidecite{Ouyang et al. (2022)}
\end{frame}



% %----------------------------------------------------------------------
% \begin{frame}{Attention and Interacting Particle System}

% Interacting Particle System: given $N$ particles $x_1,\ldots,x_N\in\R^d$,
% \[
% \dot x_i= \frac1N\sum_{j=1}^N b(x_i,x_j) = b[\mu^N_X](x_i),\qquad i=1,\ldots,N, 
% \] 
% where $\mu_X^N=N^{-1}\sum_{j=1}^N\delta_{x_j}$. 

% For one unmasked head attention with $N$ tokens, the output is
% \[
%  \operatorname{Att}_\theta(X)_i = \frac{
% \sum_{j=1}^N e^{\langle Qx_i,Kx_j\rangle/\sqrt{d_k}}Vx_j
%  }{ 
% \sum_{j=1}^N e^{\langle Qx_i,Kx_j\rangle/\sqrt{d_k}}
%  } = \Phi[\mu_X^N](x_i).
% \]


% \begin{itemize}
%   \item Finite-token: $N$ tokens in the sequence.
%   \item Mean-field: $N\to\infty$ with law $\mu_X^N\to\mu$.
%   \item Continuous depth: $L\to\infty$ with $\ell/L\to t\in[0,1]$. 
% \end{itemize}
% \end{frame}







%======================================================================
\begin{frame}{Reading and References I: Language and ICL}
\footnotesize
\begin{thebibliography}{9}
\bibitem{Vaswani2017}
A.~Vaswani et al. (2017),
``Attention Is All You Need,'' \emph{NeurIPS 30}.
\href{https://proceedings.neurips.cc/paper_files/paper/2017/hash/3f5ee243547dee91fbd053c1c4a845aa-Abstract.html}{Proceedings link}.

\bibitem{Brown2020}
T.~B.~Brown et al. (2020),
``Language Models Are Few-Shot Learners,'' \emph{NeurIPS 33}, 1877--1901.
\href{https://proceedings.neurips.cc/paper/2020/hash/1457c0d6bfcb4967418bfb8ac142f64a-Abstract.html}{Proceedings link}.

\bibitem{Xie2022}
S.~M.~Xie, A.~Raghunathan, P.~Liang, and T.~Ma (2022),
``An Explanation of In-Context Learning as Implicit Bayesian Inference,'' \emph{ICLR}.
\href{https://openreview.net/forum?id=RdJVFCHjUMI}{OpenReview link}.

\bibitem{Ouyang2022}
L.~Ouyang et al. (2022),
``Training Language Models to Follow Instructions with Human Feedback,'' \emph{NeurIPS 35}.
\href{https://arxiv.org/abs/2203.02155}{arXiv:2203.02155}.

% \bibitem{Sennrich2016}
% R.~Sennrich, B.~Haddow, and A.~Birch (2016),
% ``Neural Machine Translation of Rare Words with Subword Units,'' \emph{ACL}.
% \href{https://aclanthology.org/P16-1162/}{ACL Anthology link}.

\end{thebibliography}
 \end{frame}

%----------------------------------------------------------------------
\begin{frame}{Reading and References II: Scientific and World Models}
\scriptsize
\begin{thebibliography}{9}
\bibitem{ICON}
L.~Yang, S.~Liu, T.~Meng, and S.~J.~Osher (2023),
``In-Context Operator Learning with Data Prompts for Differential Equation Problems,'' \emph{PNAS} 120(39).
\href{https://doi.org/10.1073/pnas.2310142120}{doi:10.1073/pnas.2310142120}.
\bibitem{PDEformer}
Z.~Ye et al. (2024), ``PDEformer-1: A Foundation Model for One-Dimensional Partial Differential Equations,''
\href{https://arxiv.org/abs/2407.06664}{arXiv:2407.06664}.
\bibitem{PROSE}
Y.~Liu, Z.~Zhang, and H.~Schaeffer (2024), ``PROSE,'' \emph{Neural Networks} 180, 106707.
\href{https://doi.org/10.1016/j.neunet.2024.106707}{doi link}.
\bibitem{Poseidon}
M.~Herde et al. (2024), ``Poseidon: Efficient Foundation Models for PDEs,'' \emph{NeurIPS 37}.
\href{https://proceedings.neurips.cc/paper_files/paper/2024/hash/84e1b1ec17bb11c57234e96433022a9a-Abstract-Conference.html}{Proceedings link}.
\bibitem{WorldModels}
D.~Ha and J.~Schmidhuber (2018), ``Recurrent World Models Facilitate Policy Evolution,'' \emph{NeurIPS 31}.
\href{https://papers.nips.cc/paper_files/paper/2018/hash/2de5d16682c3c35007e4e92982f1a2ba-Abstract.html}{Proceedings link}.
\bibitem{PlanetDreamer}
D.~Hafner et al. (2019), ``Learning Latent Dynamics for Planning from Pixels,''
\href{https://proceedings.mlr.press/v97/hafner19a.html}{\emph{ICML}};
D.~Hafner et al. (2020), ``Dream to Control,''
\href{https://openreview.net/forum?id=S1lOTC4tDS}{\emph{ICLR}}.
\bibitem{MuZeroGenie}
J.~Schrittwieser et al. (2020), ``Mastering Atari, Go, Chess and Shogi by Planning with a Learned Model,''
\href{https://doi.org/10.1038/s41586-020-03051-4}{\emph{Nature} 588};
J.~Bruce et al. (2024), ``Genie: Generative Interactive Environments,''
\href{https://proceedings.mlr.press/v235/bruce24a.html}{\emph{ICML}}.
\end{thebibliography}

\end{frame}

\end{document}
